This thesis investigated whether predictive multiplicity---the disagreement among near-optimal models---exhibits spatial structure in feature space, what drives that structure, and whether it remains stable under reasonable analytical variations. To answer these questions, it introduced a framework that combines Rashomon-set construction, observation-wise multiplicity metrics, and spatial autocorrelation analysis through Moran's $I$ and LISA. The resulting perspective moves beyond aggregate disagreement summaries and instead treats multiplicity as a localized property of the data-model interaction.

\subsection{Summary of Findings}

Across the benchmark datasets and model families considered in this thesis, five main findings emerge.

\begin{enumerate}
    \item \textbf{Predictive multiplicity is spatially structured.} Observation-wise disagreement is not distributed uniformly across the test set. Instead, high-variance observations tend to cluster with other high-variance observations, producing statistically significant local hotspot regions in feature space. This means that predictive multiplicity is not only a global property of a model class, but also a localized property of specific subpopulations.

    \item \textbf{Hotspots are more informative than global averages alone.} Global disagreement metrics remain useful, but they do not identify \emph{where} instability occurs. The spatial analysis shows that two observations can be associated with similarly confident predictions while differing substantially in multiplicity exposure: one may lie in a stable region where near-optimal models agree, while the other lies in a hotspot where equally accurate models disagree strongly. Localizing multiplicity therefore changes the practical interpretation of predictive reliability.

    \item \textbf{Model family is the dominant driver, with hyperparameters adding further structure.} Between-family differences explain a substantial share of predictive variance, and this share typically increases inside hotspots. Within families, a relatively small number of hyperparameters account for much of the remaining disagreement. Multiplicity is therefore not random noise induced by training instability alone, but is systematically shaped by concrete modeling choices.

    \item \textbf{The main spatial findings are robust.} The detected hotspot structure persists under changes to the Rashomon set size $K$, the neighborhood size $k$, post-hoc calibration, and alternative graph constructions based on PCA or cosine distance. These robustness checks do not imply that every numerical value is invariant, but they do show that the core substantive conclusion---that multiplicity is spatially organized---does not depend on a single arbitrary specification.

    \item \textbf{High-multiplicity regions can be interpreted.} The rule-based analyses show that hotspots are not arbitrary statistical artifacts. They can be described by compact feature conditions with meaningful purity and lift, especially on COMPAS, where age and prior criminal history recur as important descriptors. Synthetic validation further confirms that the pipeline recovers known ambiguous regions with high precision and recall, which strengthens confidence that the real-data hotspots reflect genuine structure.
\end{enumerate}

Taken together, these results support the central claim of the thesis: predictive multiplicity is not merely a global nuisance metric, but a structured and localized phenomenon that can be detected, quantified, and described.

\subsection{Interpretation}

The broader implication of these findings is that not all predictions produced by an accurate model should be treated as equally reliable. Standard evaluation focuses on whether a selected model performs well on average. The results of this thesis show that average performance is only part of the picture. Even when a model is chosen from a set of near-optimal alternatives, its predictions may be much less stable for some individuals than for others because equally good models would have produced materially different outputs.

This perspective also clarifies the difference between predictive multiplicity and more familiar notions of uncertainty. Classical uncertainty analyses often ask whether a model is uncertain because the data are noisy, the sample is limited, or the fitted parameters vary. Rashomon multiplicity addresses a different source of unreliability: the existence of many acceptable models that disagree in behavior. The contribution of this thesis is to show that this form of instability can itself have spatial structure. In other words, model choice matters more in some parts of the feature space than in others.

From an applied perspective, this suggests a different deployment logic. Rather than treating all predictions from the selected model as equally trustworthy, practitioners can distinguish between stable regions and multiplicity hotspots. Predictions in hotspot regions may warrant caution, additional review, or abstention, especially in high-stakes settings such as criminal justice or credit scoring. The contribution of the thesis is therefore not only descriptive, but also diagnostic: it identifies where arbitrary model choice is most consequential.

At the same time, the results highlight an important trade-off. Restricting the hypothesis space to a narrower family may reduce multiplicity, but it also removes part of the diversity that makes Rashomon-set analysis informative in the first place. A system designer must therefore decide whether the goal is to minimize disagreement, to characterize it, or to make it visible for downstream governance decisions.

\subsection{Limitations}

Several limitations should be acknowledged.

First, the Rashomon set is approximated through a finite candidate pool and top-$K$ selection. This is a practical and reproducible operationalization, but it cannot fully capture the potentially continuous set of all near-optimal models. The observed multiplicity structure therefore depends on the chosen model families and hyperparameter search spaces.

Second, the spatial analysis depends on the graph used to represent local structure in feature space. Although the main conclusions are robust across several alternatives, the exact size and shape of hotspot regions still reflect design choices in graph construction and distance representation.

Third, the local hotspot assignments are more stable at the region level than at the exact point level. The same broader parts of feature space tend to recur across runs, but the precise HH membership can shift. The method is therefore better interpreted as identifying unstable \emph{areas} rather than assigning an immutable hotspot label to each individual observation.

Fourth, the fairness analysis is descriptive rather than causal. Differences in subgroup exposure indicate that some groups may be more affected by multiplicity, but these results do not establish why such disparities arise.

Finally, the empirical study is based on a limited number of benchmark datasets and binary classification tasks. Although the pattern is consistent across the settings studied here, broader validation on additional domains would strengthen the generality of the conclusions.

\subsection{Future Work}

Several extensions follow naturally from this work.

\begin{itemize}
    \item \textbf{Continuous Rashomon approximations.} Replacing the discrete model pool with continuous or optimization-based approximations could provide a tighter view of multiplicity and reduce dependence on sampled candidate sets.

    \item \textbf{Multiplicity-aware deployment policies.} Hotspot information could be incorporated into selective prediction, abstention, or human-in-the-loop review workflows, so that especially model-contingent cases are handled differently from stable ones.

    \item \textbf{Temporal analysis under distribution shift.} Studying whether hotspot locations remain stable over time would connect multiplicity analysis to drift monitoring and deployment reliability.

    \item \textbf{Stronger fairness analysis.} Future work could move from descriptive subgroup exposure to hierarchical or causal analyses that investigate why some groups are more exposed to multiplicity than others.

    \item \textbf{Extensions beyond binary classification.} The framework could be generalized to multiclass classification, regression, or ranking settings, where disagreement may take different geometric forms.
\end{itemize}

In summary, this thesis shows that a spatial perspective reveals structure in predictive multiplicity that aggregate metrics alone cannot capture. By localizing disagreement, identifying its drivers, and characterizing the affected regions, the proposed framework provides a foundation for treating multiplicity as a practical reliability and governance issue rather than as a purely abstract property of model classes.